%==============================================================
% Advanced Scientific Writing -- Assignment 5
% Research-paper SECOND DRAFT, typeset with the official PSB / World
% Scientific author kit (ws-procs11x85 document class).
%
% Build:  pdflatex twice  (or: latexmk -pdf -cd)
%==============================================================

\documentclass{ws-procs11x85}
\usepackage{ws-procs-thm}
\usepackage{booktabs}

\begin{document}

\title{Nicotinic Acetylcholine Receptor Prediction using VEP for Humans}

\author{Hari Krishna Mohan Raj}

\address{Department of Computational Chemistry,\\
Hochschule Bonn Rhein Sieg, Sankt Augustin, Germany\\
E-mail: harikrishnam0711@gmail.com}

\begin{abstract}
Missense variants in human nicotinic acetylcholine receptor (nAChR)
genes are linked to congenital myasthenic syndromes, epilepsy, and
related disorders, yet most newly observed variants are of uncertain
significance. Standard variant effect predictors report only whether a
variant is damaging, not whether it causes loss of function (LOF) or
gain of function (GOF), a distinction that is decisive for therapy. We
assemble a curated dataset of experimentally characterised nAChR
missense variants labelled by direction of effect, engineer sequence-
and structure-derived features, and benchmark a panel of machine-learning
classifiers to predict LOF versus GOF. Gradient boosted models perform
best, indicating that direction-aware, receptor-specific prediction is
feasible and can help triage nAChR variants of uncertain significance
toward functional follow-up and precision therapy.
\end{abstract}

\keywords{Variant effect prediction; nicotinic acetylcholine receptor;
gain of function; loss of function; ion channel; machine learning.}

% required, do-not-remove
\copyrightinfo{\copyright\ 2026 Hari Krishna Mohan Raj, Karl Kirschner. Open Access chapter
published by World Scientific Publishing Company and distributed under
the terms of the Creative Commons Attribution Non-Commercial (CC
BY-NC) 4.0 License.}

%==============================================================
\section{Introduction}\label{sec:intro}
%==============================================================
Missense variants in the genes encoding nicotinic acetylcholine
receptors (nAChRs) underlie congenital myasthenic syndromes at the
neuromuscular junction and several inherited epilepsies in the brain.
As clinical sequencing becomes routine, new variants are observed far
faster than their functional consequences can be measured, leaving most
as variants of uncertain significance (VUS). Existing computational
predictors such as SIFT, PolyPhen-2,
CADD, and AlphaMissense report whether a
variant is damaging, but not whether it causes loss of function (LOF) or
gain of function (GOF). For nAChRs this is a critical omission, because
LOF and GOF often demand opposite clinical strategies. This work asks
whether the direction of effect can be predicted computationally, and
treats it as a receptor-specific binary classification problem.

%==============================================================
\section{Background}\label{sec:background}
%==============================================================
nAChRs are pentameric ligand-gated ion channels assembled from a
repertoire of seventeen human subunit genes that combine into muscle-type
and neuronal receptors. The affected subunit largely determines where a
receptor acts and which disorder a variant is associated with, so
subunit identity is as informative as the substitution itself. The
direction of a variant's effect is measurable by electrophysiology, which
reveals whether a mutant receptor conducts too little or too much current,
but such assays are slow and costly and cannot scale to the volume of
variants now being reported. This gap between a clinically decisive
property and the cost of measuring it directly motivates a computational
approach.

\subsection{Existing Computational Approaches and Research Gaps}

Several tools attempt to predict variant pathogenicity or direction of
effect, but none are tailored to nAChR biology.
Table~\ref{tab:tools} summarises the key distinctions.
General pathogenicity tools such as SIFT~\cite{sift},
PolyPhen-2~\cite{polyphen2}, CADD~\cite{cadd}, and
AlphaMissense~\cite{alphamissense} output a damage score but provide no
information on whether a variant causes LOF or GOF.
Direction-aware tools exist but operate at a different scope:
funNCion~\cite{funncion} is trained on sodium and calcium channel
variants (ROC~$\approx$~0.85) and does not generalise to nAChRs;
LoGoFunc~\cite{logofunc} predicts genome-wide LOF/GOF/neutral but
treats all gene families uniformly; PreMode~\cite{premode} uses a graph
neural network for mode-of-action prediction without channel-family
specificity; and MissION~\cite{mission}, the most recent tool
(ROC-AUC~0.925 across 47 ion-channel genes), uses protein language
model embeddings and structural features but aggregates across a broad
set of ion channels without nAChR-specific subunit modelling.
The gap this work addresses is a direction-aware predictor built
specifically around nAChR subunit identity and the conformational
context that distinguishes agonist-bound open from resting closed states.

\begin{table}[ht]
\tbl{Comparison of computational approaches for variant effect
prediction in the context of this work.}
{\begin{tabular}{@{}lcccc@{}}
\toprule
Tool & LOF/GOF & Channel-specific & nAChR-specific \\
\midrule
SIFT                  & No  & No          & No  \\
AlphaMissense         & No  & No          & No  \\
funNCion              & Yes & Na/Ca only  & No  \\
LoGoFunc              & Yes & No          & No  \\
PreMode               & Yes & No          & No  \\
MissION               & Yes & 47 channels & No  \\
\textbf{This work}    & \textbf{Yes} & \textbf{nAChR} & \textbf{Yes} \\
\bottomrule
\end{tabular}}
\label{tab:tools}
\end{table}

%==============================================================
\section{Analytical Objective}\label{sec:objective}
%==============================================================
The objective is to classify a given nAChR missense variant as LOF or
GOF, rather than merely pathogenic or benign. Unlike general pathogenicity
tools, and unlike direction-aware predictors developed for other channel
families such as funNCion, LoGoFunc, and
PreMode, the model is built specifically around nAChR
biology and subunit identity. The task is framed as supervised binary
classification, evaluated by its ability to recover the experimentally
established direction of effect for held-out variants.

%==============================================================
\section{Experimental Data}\label{sec:data}
%==============================================================
The training data comprise a curated set of experimentally characterised
nAChR missense variants spanning the majority of human subunit genes,
each labelled LOF or GOF from primary electrophysiology
reports. The current dataset contains approximately 350 variants
(218 LOF, 133 GOF) across twelve subunit genes. Every variant is
represented by an engineered feature vector combining amino-acid
physicochemical properties, evolutionary substitution scores (BLOSUM62
and Grantham distance), normalised sequence position, a one-hot encoding
of subunit identity, and structural descriptors derived from experimental
cryo-EM structures (per-residue B-factor, secondary structure, and local
packing density). A majority of variants map successfully onto solved
structures, providing the structural context that sequence-only methods
lack.

\subsection{Preliminary Evidence from a Related Ion Channel}

The feasibility of direction-aware prediction for ion channels is
supported by a parallel VEP project applied to epithelial sodium channel
(ENaC) variants, which serves as the template for the present work.
In ENaC, LOF variants cause pseudohypoaldosteronism type~1 while GOF
variants underlie Liddle syndrome, producing opposite clinical
phenotypes that demand opposite treatments, mirroring the LOF/GOF
distinction in nAChR. The ENaC project demonstrated that sequence and
structural features are sufficient to separate these two directions of
effect above chance, providing preliminary evidence that
receptor-specific direction prediction is achievable for pentameric ion
channels without large training sets~\cite{enac}.

%==============================================================
\section{Features, Algorithm and Parameters}\label{sec:algorithms}
%==============================================================
A panel of seven machine learning classifiers was benchmarked on the same
data and features, ranging from linear models and support vector machines
to random forests, gradient-boosted trees (XGBoost and LightGBM), and a
neural network. Models were trained and tuned with nested cross-validation
and Optuna-based hyperparameter optimisation to obtain unbiased estimates
on this small dataset. Consistent with expectations for limited
tabular data, the gradient-boosted tree methods outperformed the linear,
kernel, and neural alternatives.

There are approximately 50 features in total, organised into four
categories as summarised in Table~\ref{tab:features}. Ablation studies
to quantify the contribution of each category remain ongoing.

\begin{table}[ht]
\tbl{Overview of engineered feature categories used in the classifier.}
{\begin{tabular}{@{}ll@{}}
\toprule
Category & Features \\
\midrule
Sequence / physicochemical &
  Hydrophobicity, charge, size, polarity,\\
  & aromaticity of wild-type and mutant residue;\\
  & change in each property ($\Delta$) \\
Evolutionary &
  BLOSUM62 substitution score, Grantham distance \\
Positional &
  Normalised position along the subunit sequence \\
Subunit identity &
  One-hot encoding of the 17 nAChR subunit genes \\
Structural &
  Per-residue B-factor, secondary structure class\\
  & (helix / strand / coil), local packing density \\
\bottomrule
\end{tabular}}
\label{tab:features}
\end{table}

%==============================================================
\section{Metrics}\label{sec:metrics}
%==============================================================
Performance was assessed with five-fold cross-validation, reporting
accuracy and the F1 score for GOF detection given the class imbalance.
The strongest models reached an F1 of roughly 0.67 at about 75\% accuracy,
well above a random baseline near 0.38, indicating genuine signal in the
engineered features. Adding the structural descriptors produced a modest
but consistent improvement over sequence-only features.

[TODO: add a compact results table and confidence intervals once the
final feature set and dataset are locked.]

%==============================================================
\section{Conclusion}\label{sec:conclusion}
%==============================================================
These results indicate that direction-aware, receptor-specific prediction
of nAChR variant effects is feasible: engineered sequence and structural
features carry enough signal to separate LOF from GOF above chance. The
approach fills a clear gap in the existing landscape, where no current
tool combines direction-of-effect prediction with nAChR subunit-level
specificity. Such a predictor could triage the expanding population of
nAChR VUS, prioritising which variants warrant costly functional
follow-up and supporting precision therapy. The main limitations are data
scarcity and the treatment of LOF and GOF as a clean dichotomy; future
work should expand the curated dataset, establish a shared
direction-of-effect benchmark for nAChRs, and explore whether protein
language model embeddings can further improve performance on this
limited training set.

%==============================================================
% References
%==============================================================
\begin{thebibliography}{9}

\bibitem{sift} P.~C. Ng and S.~Henikoff, SIFT: predicting amino acid
changes that affect protein function, {\em Nucleic Acids Res.}
{\bf 31}, 3812 (2003).

\bibitem{polyphen2} I.~A. Adzhubei {\em et al.}, A method and server
for predicting damaging missense mutations, {\em Nat. Methods}
{\bf 7}, 248 (2010).

\bibitem{cadd} M.~Kircher {\em et al.}, A general framework for
estimating the relative pathogenicity of human genetic variants,
{\em Nat. Genet.} {\bf 46}, 310 (2014).

\bibitem{alphamissense} J.~Cheng {\em et al.}, Accurate proteome-wide
missense variant effect prediction with AlphaMissense, {\em Science}
{\bf 381}, eadg7492 (2023).

\bibitem{esm1b} N.~Brandes {\em et al.}, Genome-wide prediction of
disease variant effects with a deep protein language model,
{\em Nat. Genet.} {\bf 55}, 1512 (2023).

\bibitem{funncion} A.~Brunklaus {\em et al.}, Gene variant effects
across sodium channelopathies predict function and guide precision
therapy, {\em Brain} {\bf 145}, 4275 (2022).

\bibitem{logofunc} D.~Stein {\em et al.}, Genome-wide prediction of
pathogenic gain- and loss-of-function variants from ensemble learning
of a diverse feature set, {\em Genome Med.} {\bf 15}, 103 (2023).

\bibitem{premode} G.~Zhang {\em et al.}, Predicting the mode of action
of missense variants (PreMode), bioRxiv (2024).

\bibitem{mission} [Author names TBD], MissION: direction-of-effect
prediction for ion-channel missense variants via protein language model
embeddings and structural features, {\em medRxiv} (2025).

\bibitem{enac} R.~P. Lifton {\em et al.}, Molecular mechanisms of
human hypertension, {\em Cell} {\bf 104}, 545 (2001).

\bibitem{mohanraj2026} H.~K. Mohan Raj {\em et al.}, A curated dataset of
experimentally characterised missense variants in human nicotinic
acetylcholine receptors, (2026).

\end{thebibliography}

\end{document}
